% The fork rule: one function, two versions, and who decides. Prints as
% Figure 4, single column, beside Section 4.3. Same panels, glyphs and palette
% as Figures 1 to 3. The instance is the one the walkthrough reaches in 5.5.
\newcommand{\FigForkCaption}{Two people change \protect\fsym{api.py::handle}
from the same saved version, one cutting three lines from the top of it and one
adding a line forty lines below. Git merges the two and prints nothing, and
nobody has run the function in the state the merge produces. \sgt{} records both
versions, includes neither, and names the function. Of everything \sgt{} refuses,
only this refusal can stop work that would have been correct.}

\newcommand{\FigForkDescription}{Three stacked panels. The top panel shows one
recorded version of a function splitting into two, one change removing three
lines at the top of the function and one adding a line at the bottom, each
marked with a small figure of a person, with the two resulting function bodies
drawn side by side. The middle panel shows git's result, a single function body
with the top lines gone and the bottom line added, annotated as a state nobody
ran. The bottom panel shows sgt's result, both versions drawn as dotted outlines
to mean neither is included, and one solid version after the developer writes
the merge and the test command passes.}

\begin{tikzpicture}[x=1mm, y=1mm, line cap=round, font=\sffamily]

% one function body. #1 x, #2 y of top edge, #3 mode: t / b / a
% t = the three lines at the top are gone.  b = a line added at the bottom.
\newcommand{\fbody}[3]{%
  \draw[card, color=sgtGrey!45, fill=white] (#1,{#2-9.4}) rectangle ({#1+13},#2);
  \foreach \yy in {-1.4,-2.6,-3.8,-5.0,-6.2,-7.4,-8.6}
    {\fill[sgtGrey!25] ({#1+1.3},{#2+\yy-0.3}) rectangle ({#1+9.4},{#2+\yy+0.15});}
  \if#3b\else
    \fill[white] ({#1+0.9},{#2-4.4}) rectangle ({#1+12.1},{#2-0.9});
    \draw[fCache, line width=0.3pt, densely dotted]
      ({#1+1.3},{#2-4.3}) rectangle ({#1+11.2},{#2-1.0});
  \fi
  \if#3t\else
    \fill[fIdx] ({#1+1.3},{#2-8.9}) rectangle ({#1+11.2},{#2-8.45});
  \fi}
% #1 x, #2 y of body top, #3 colour, #4 letter
\newcommand{\fbodytag}[4]{%
  \node[tag, fill=#3, text=white, anchor=south west,
        font=\sffamily\bfseries\fontsize{5.0}{5.8}\selectfont] at (#1,{#2+0.5}) {#4};}

% ============== 1: where the two versions come from ==================
\panel{0}{38}{84}{24}{sgtInk}{1\ \ ONE SAVED VERSION, TWO NEW ONES}

\draw[rail=sgtGrey] (2.5,48.5) -- (9.0,48.5);
\opext{5.5}{48.5}{sgtGrey}
\draw[raildim] (9.0,48.5) -- (13.5,53.0);
\draw[raildim] (9.0,48.5) -- (13.5,44.0);
\draw[rail=fCache] (13.5,53.0) -- (20.5,53.0); \oprew{17.0}{53.0}{fCache}
\draw[rail=fIdx]   (13.5,44.0) -- (20.5,44.0); \opext{17.0}{44.0}{fIdx}
\whodev{23.2}{53.0}{fCache}
\whodev{23.2}{44.0}{fIdx}

\node[anchor=west, text width=21mm, align=left, inner sep=0, text=fCache,
      font=\sffamily\fontsize{5.4}{6.2}\selectfont] at (26.4,53.2)
  {\textbf{A} the cache revert cut 3 lines off the top};
\node[anchor=west, text width=21mm, align=left, inner sep=0, text=fIdx,
      font=\sffamily\fontsize{5.4}{6.2}\selectfont] at (26.4,44.2)
  {\textbf{B} a colleague logged a request id at the end};

\fbody{53.0}{57.6}{t}\fbodytag{53.0}{57.6}{fCache}{A}
\fbody{68.5}{57.6}{b}\fbodytag{68.5}{57.6}{fIdx}{B}
\node[acc=sgtInk!70, anchor=north] at (67.2,46.8) {forty lines apart};

% ===================== 2: what git produces ==========================
\panel{0}{18}{84}{18}{sgtGrey}{2\ \ WHAT GIT PRODUCES}

\node[anchor=north west, text width=46mm, align=left, inner sep=0, text=sgtInk!85,
      font=\sffamily\fontsize{5.7}{6.7}\selectfont] at (2.2,29.8)
  {The two changes touch no common line, so git merges them and prints nothing.
   Nobody has run this function without its cache lookup and with the new log
   line.};

\fbody{53.0}{30.2}{a}
\fbodytag{53.0}{30.2}{sgtGrey}{A+B}
\annot{69.5}{24.4}{66.6}{25.4}{18}{fExport}
\node[anchor=west, text width=13mm, align=left, inner sep=0, text=fExport,
      font=\sffamily\bfseries\fontsize{5.4}{6.2}\selectfont] at (70.0,24.4)
  {a state nobody ran};

% ===================== 3: what sgt produces ==========================
\panel{0}{0}{84}{16}{sgtInk}{3\ \ WHAT \textsf{sgt} PRODUCES}

\node[anchor=north west, inner sep=0, text=sgtInk,
      font=\ttfamily\fontsize{5.6}{6.4}\selectfont] at (2.2,11.6)
  {\$ sgt resolve api.py::handle};

\draw[raildim] (2.4,5.6) -- (6.0,7.8);
\draw[raildim] (2.4,5.6) -- (6.0,3.4);
\draw[raildim] (6.0,7.8) -- (11.6,7.8); \opghost{8.8}{7.8}{fCache}
\draw[raildim] (6.0,3.4) -- (11.6,3.4); \opghost{8.8}{3.4}{fIdx}
\node[anchor=west, text width=20mm, align=left, inner sep=0, text=sgtInk!80,
      font=\sffamily\fontsize{5.4}{6.2}\selectfont] at (13.4,5.6)
  {both recorded,\\neither included};

\draw[-{Straight Barb[length=0.9mm,width=1.0mm]}, line width=0.35pt, color=sgtGrey]
  (35.0,5.6) -- (38.6,5.6);
\draw[rail=sgtInk] (40.6,5.6) -- (47.6,5.6); \oprew{44.1}{5.6}{sgtInk!70}
\whodev{50.4}{5.6}{sgtInk!70}
\node[anchor=west, text width=29mm, align=left, inner sep=0, text=sgtInk!85,
      font=\sffamily\fontsize{5.4}{6.2}\selectfont] at (53.0,5.6)
  {the version the developer wrote, included once the tests pass};

\end{tikzpicture}
